\section*{अध्याय \ref{ch:b2:current-potential-curves} --- धारा--विभव वक्र}

\begin{solution}{exo:b2:current-potential-curves:1}
ज़िंक ऑक्सीकृत होता है: \omterm{def:b2:current-potential-curves:ie-curve}{ऐनोडी धारा}, धनात्मक। ताँबे के आयन ताँबे के इलेक्ट्रोड पर
अपचित होते हैं: \omterm{def:b2:current-potential-curves:ie-curve}{कैथोडी धारा}, ऋणात्मक। दोनों धाराएँ परिमाण में
बराबर हैं, क्योंकि परिपथ से वही आवेश बहता है।
\end{solution}

\begin{solution}{exo:b2:current-potential-curves:2}
$|i_{\lim}| = 96\,485 \times 0.50 \times 10^{-4} \times 1.6 \times 10^{-9} \times
2.0/(25 \times 10^{-6}) = \qty{6.2e-4}{A}$, अर्थात् \qty{0.62}{mA}।
\end{solution}

\begin{solution}{exo:b2:current-potential-curves:3}
\omterm{def:b2:current-potential-curves:fast-slow}{तीव्र निकाय} अपने नेर्न्स्ट विभव पर विभव अक्ष को तीव्रता से काटता है; \omterm{def:b2:current-potential-curves:fast-slow}{मंद
निकाय} उसके आसपास शून्य धारा का एक समतल खंड दिखाता है, उसकी शाखाएँ
\omterm{def:b2:current-potential-curves:overpotential}{अधिविभव} के बाद ही आरंभ होती हैं। प्लैटिनम पर: \ce{Fe^3+}/\ce{Fe^2+} (तीव्र);
\ce{O2}/\ce{H2O} (मंद)।
\end{solution}

\begin{solution}{exo:b2:current-potential-curves:4}
अर्ध-तरंग विभव पर, जो $E^\circ = \qty{0.77}{V}$ के बराबर है जब दोनों रूप
समान रूप से विसरित होते हैं।
\end{solution}

\begin{solution}{exo:b2:current-potential-curves:5}
लगभग \qty{0.77}{V} पर \ce{Fe^3+} की एक कैथोडी तरंग (तीव्र), जो उसकी सांद्रता के
समानुपाती पठार तक पहुँचती है। लगभग \qty{0.34}{V} पर दूसरी तरंग, ताँबा
जमता हुआ; उसका पठार पहले में जुड़ता है, समान
सांद्रता और समान विसरण गुणांकों के लिए लगभग दुगुना ऊँचा, क्योंकि $n = 2$। फिर,
प्लैटिनम पर pH 0 पर लगभग \qty{0}{V} पर (इलेक्ट्रोड के ताँबे से लेपित हो जाने पर नीचे),
\ce{H+} के अपचयन की खड़ी भित्ति।
\end{solution}

\begin{solution}{exo:b2:current-potential-curves:6}
\qty{0.50}{V} पर: केवल \ce{Fe^3+} का \ce{Fe^2+} में अपचयन होता है, अपनी \omterm{def:b2:current-potential-curves:limiting-current}{सीमांत धारा} पर।
\qty{0.10}{V} पर: \ce{Fe^3+} अपचित होता है और साथ ही ताँबा जमता है,
हर एक अपनी \omterm{def:b2:current-potential-curves:limiting-current}{सीमांत धारा} पर; कुल धारा दोनों पठारों का योग है।
\end{solution}

\begin{solution}{exo:b2:current-potential-curves:7}
pH 0: \qty{0}{V} और \qty{1.23}{V}; pH 14: \qty{-0.83}{V} और \qty{0.40}{V}। प्लैटिनम
पर खिड़की अपनी चौड़ाई बनाए रखती है और प्रति pH इकाई \qty{59}{mV} नीचे खिसकती है,
हर भित्ति अपने युग्म का अनुसरण करती है, ऐनोडी भित्ति अब भी ऑक्सीजन
\omterm{def:b2:current-potential-curves:overpotential}{अधिविभव} से ऊपर खिसकी रहती है।
\end{solution}

\begin{solution}{exo:b2:current-potential-curves:8}
\qty{1.0}{V} पर \ce{Fe^2+} अपने पठार पर ऑक्सीकृत होता है (\omterm{def:b2:current-potential-curves:ie-curve}{ऐनोडी धारा}
$[\ce{Fe^2+}]$ के समानुपाती) और कोई भी \ce{Ce^4+} अपने पठार पर अपचित होता है
(\omterm{def:b2:current-potential-curves:ie-curve}{कैथोडी धारा} $[\ce{Ce^4+}]$ के समानुपाती)। तुल्यता बिंदु से पहले
\omterm{def:b2:current-potential-curves:ie-curve}{ऐनोडी धारा} मिलाए गए आयतन के साथ रैखिक रूप से गिरती है; उसके बाद \omterm{def:b2:current-potential-curves:ie-curve}{कैथोडी
धारा} रैखिक रूप से बढ़ती है। दोनों सरल रेखाएँ तुल्यता बिंदु पर शून्य को काटती हैं।
\end{solution}

\begin{solution}{exo:b2:current-potential-curves:9}
पठार दुगुने हो जाते हैं ($|i_{\lim}| \propto 1/\delta$)। अर्ध-तरंग विभव
नहीं बदलता, क्योंकि परत दोनों रूपों के लिए समान है; \omterm{def:b2:current-potential-curves:fast-slow}{तीव्र निकाय} का शून्य-धारा
विभव भी नहीं, जो नेर्न्स्ट विभव है।
\end{solution}

\begin{solution}{exo:b2:current-potential-curves:10}
अध्याय की तरह, $i_{\lim,c} = 0$ के साथ: $E = E_{1/2} + (RT/nF)\ln[i/(i_{\lim} -
i)]$। दशमलव लघुगणकों में $E = E_{1/2} + (RT\ln 10/nF)\log[i/(i_{\lim} - i)]$,
जो \qty{298}{K} पर $0.0592/n$ वोल्ट प्रवणता वाली सरल रेखा है, और
$E$ अक्ष को $E_{1/2}$ पर काटती है।
\end{solution}

\begin{solution}{exo:b2:current-potential-curves:11}
अकेला ज़िंक उस विभव पर स्थिर होता है जहाँ उसकी \omterm{def:b2:current-potential-curves:ie-curve}{ऐनोडी धारा} (ज़िंक का तीव्र
ऑक्सीकरण) ज़िंक पर \ce{H+} के अपचयन की \omterm{def:b2:current-potential-curves:ie-curve}{कैथोडी धारा} के बराबर होती है,
जो छोटी है क्योंकि वह अपचयन ज़िंक पर मंद है: विलयन
मंद है। प्लैटिनम को छूने पर ज़िंक द्वारा मुक्त इलेक्ट्रॉन प्लैटिनम पर \ce{H+} का
अपचयन कर सकते हैं, जहाँ अपचयन तीव्र है: दिए गए विभव पर \omterm{def:b2:current-potential-curves:ie-curve}{कैथोडी धारा}
कहीं बड़ी होती है, संतुलन ऊँची धारा पर पहुँचता है, और
ज़िंक तेज़ी से घुलता है जबकि प्लैटिनम पर हाइड्रोजन बनती है।
\end{solution}

\begin{solution}{exo:b2:current-potential-curves:12}
प्लैटिनम पर जल का ऑक्सीकरण pH 0 पर लगभग \qty{1.6}{V} पर आरंभ होता है:
\qty{2.0}{V} से पहले धारा डाइऑक्सीजन बनाने में खर्च होती। ऐसा इलेक्ट्रोड जिस पर
जल के ऑक्सीकरण को बड़ा \omterm{def:b2:current-potential-curves:overpotential}{अधिविभव} चाहिए, भित्ति को
\qty{2.0}{V} से आगे धकेल देता है, जिससे स्पीशीज़ खिड़की के भीतर ऑक्सीकृत हो सके।
\end{solution}

\begin{solution}{pb:b2:current-potential-curves:1}
\textbf{1.} \ce{O2 + 2H2O + 4e- -> 4OH-}।
\textbf{2.} \ce{Ag + Cl- -> AgCl + e-} (चार बार); समग्र \ce{O2 + 2H2O + 4Ag +
4Cl- -> 4AgCl + 4OH-}।
\textbf{3.} $E = 1.229 - 0.0592 \times 7 + (0.0592/4)\log 0.21 = \qty{0.80}{V}$।
\textbf{4.} सक्रियता 1 वाले क्लोराइड के साथ $E^\circ(\ce{AgCl}/\ce{Ag}) = \qty{0.22}{V}$:
हाइड्रोजन पैमाने पर $-0.60 + 0.22 = \qty{-0.38}{V}$।
\textbf{5.} \ce{O2} का अपचयन \omterm{def:b2:current-potential-curves:fast-slow}{मंद निकाय} है: अपने नेर्न्स्ट विभव से नीचे
धारा बढ़ने से पहले वोल्ट के कई दसवें भाग का \omterm{def:b2:current-potential-curves:overpotential}{अधिविभव} चाहिए,
और पठार तक पहुँचने के लिए और अधिक।
\textbf{6.} पठार पर धारा केवल \ce{O2} की आपूर्ति पर निर्भर करती है,
विभव के छोटे विचलनों पर नहीं: यह सांद्रता मापती है।
\textbf{7.} जल का डाइहाइड्रोजन में अपचयन (कैथोडी भित्ति), जो
ऑक्सीजन से असंबंधित धारा जोड़ देता।
\textbf{8.} झिल्ली के आर-पार रैखिक, नमूने वाले पक्ष पर $c$ से कैथोड पर
शून्य तक।
\textbf{9.} प्रति इकाई क्षेत्रफल $j = D_m c/L$।
\textbf{10.} $i = -4FAD_m c/L$ (कैथोडी), $|i| = 4FAD_mc/L$।
\textbf{11.} झिल्ली सबसे मंद चरण है: उसके आर-पार विसरण बाहर के जल की तुलना में
कहीं धीमा है, इसलिए झिल्ली \omterm{def:b2:current-potential-curves:limiting-current}{विसरण परत} की भूमिका
निभाती है।
\textbf{12.} सांद्रता प्रवणता झिल्ली के भीतर है, जिसकी मोटाई
विलोडन पर निर्भर नहीं करती (बशर्ते बाहर का जल नवीकृत होता रहे);
झिल्ली पर जमाव मोटाई बढ़ाता है और धारा घटाता है।
\textbf{13.} $0.2095 \times (101.325 - 3.17) = \qty{20.56}{kPa}$।
\textbf{14.} $c = 1.3 \times 10^{-5} \times 20\,560 = \qty{0.27}{mol/m^3}$ ($0.267$)।
\textbf{15.} $0.267 \times 32.00 = \qty{8.6}{mg/L}$ (\unit{g/m^3})।
\textbf{16.} $|i| = 4 \times 96\,485 \times 1.0 \times 10^{-5} \times 1.0 \times 10^{-10}
\times 0.267/(25 \times 10^{-6}) = \qty{4.1}{\micro A}$।
\textbf{17.} प्रति \unit{mg/L} $4.05/8.55 = \qty{0.474}{\micro A}$।
\textbf{18.} झिल्ली की मोटाई और उसका विसरण गुणांक
यथार्थता से ज्ञात नहीं होते और आयु तथा ताप के साथ बदलते हैं; किसी ज्ञात
विलयन में मापन उन सबको सम्मिलित कर लेता है।
\textbf{19.} शून्य (व्यवहार में एक छोटी अवशिष्ट धारा, जिसे घटा दिया जाता है)।
\textbf{20.} $2.80/0.474 = \qty{5.9}{mg/L}$।
\textbf{21.} संतृप्ति का $2.80/4.05 = 69~\%$।
\textbf{22.} कार्बनिक पदार्थ (मल-जल, मृत पौधे) के अपघटन द्वारा उपभुक्त ऑक्सीजन,
जो वायु और प्रकाश-संश्लेषण द्वारा उसकी आपूर्ति से तेज़ी से उपभुक्त होती है।
\textbf{23.} प्रति घंटा $2.80 \times 10^{-6}/(4 \times 96\,485) \times 3600 = \qty{2.6e-8}{mol}$:
एक लीटर जल की भी ऑक्सीजन (\qty{1.8e-4}{mol}) की तुलना में
नगण्य।
\textbf{24.} \ce{O2} की विलेयता (\omterm{prop:b2:chemical-potential:henry}{हेनरी स्थिरांक}) और झिल्ली के आर-पार
विसरण, दोनों ताप के साथ बदलते हैं: \qty{25}{\degreeCelsius} पर ज्ञात
सुग्राहिता अब लागू नहीं होती।
\textbf{25.} नमूने में $\boldsymbol{\approx \qty{5.9}{mg/L}}$ घुली
ऑक्सीजन है।
\end{solution}
